%% ODER: format ==         = "\mathrel{==}"
%% ODER: format /=         = "\neq "
%
%
\makeatletter
\@ifundefined{lhs2tex.lhs2tex.sty.read}%
  {\@namedef{lhs2tex.lhs2tex.sty.read}{}%
   \newcommand\SkipToFmtEnd{}%
   \newcommand\EndFmtInput{}%
   \long\def\SkipToFmtEnd#1\EndFmtInput{}%
  }\SkipToFmtEnd

\newcommand\ReadOnlyOnce[1]{\@ifundefined{#1}{\@namedef{#1}{}}\SkipToFmtEnd}
\usepackage{amstext}
\usepackage{amssymb}
\usepackage{stmaryrd}
\DeclareFontFamily{OT1}{cmtex}{}
\DeclareFontShape{OT1}{cmtex}{m}{n}
  {<5><6><7><8>cmtex8
   <9>cmtex9
   <10><10.95><12><14.4><17.28><20.74><24.88>cmtex10}{}
\DeclareFontShape{OT1}{cmtex}{m}{it}
  {<-> ssub * cmtt/m/it}{}
\newcommand{\texfamily}{\fontfamily{cmtex}\selectfont}
\DeclareFontShape{OT1}{cmtt}{bx}{n}
  {<5><6><7><8>cmtt8
   <9>cmbtt9
   <10><10.95><12><14.4><17.28><20.74><24.88>cmbtt10}{}
\DeclareFontShape{OT1}{cmtex}{bx}{n}
  {<-> ssub * cmtt/bx/n}{}
\newcommand{\tex}[1]{\text{\texfamily#1}}	% NEU

\newcommand{\Sp}{\hskip.33334em\relax}


\newcommand{\Conid}[1]{\mathit{#1}}
\newcommand{\Varid}[1]{\mathit{#1}}
\newcommand{\anonymous}{\kern0.06em \vbox{\hrule\@width.5em}}
\newcommand{\plus}{\mathbin{+\!\!\!+}}
\newcommand{\bind}{\mathbin{>\!\!\!>\mkern-6.7mu=}}
\newcommand{\rbind}{\mathbin{=\mkern-6.7mu<\!\!\!<}}% suggested by Neil Mitchell
\newcommand{\sequ}{\mathbin{>\!\!\!>}}
\renewcommand{\leq}{\leqslant}
\renewcommand{\geq}{\geqslant}
\usepackage{polytable}

%mathindent has to be defined
\@ifundefined{mathindent}%
  {\newdimen\mathindent\mathindent\leftmargini}%
  {}%

\def\resethooks{%
  \global\let\SaveRestoreHook\empty
  \global\let\ColumnHook\empty}
\newcommand*{\savecolumns}[1][default]%
  {\g@addto@macro\SaveRestoreHook{\savecolumns[#1]}}
\newcommand*{\restorecolumns}[1][default]%
  {\g@addto@macro\SaveRestoreHook{\restorecolumns[#1]}}
\newcommand*{\aligncolumn}[2]%
  {\g@addto@macro\ColumnHook{\column{#1}{#2}}}

\resethooks

\newcommand{\onelinecommentchars}{\quad-{}- }
\newcommand{\commentbeginchars}{\enskip\{-}
\newcommand{\commentendchars}{-\}\enskip}

\newcommand{\visiblecomments}{%
  \let\onelinecomment=\onelinecommentchars
  \let\commentbegin=\commentbeginchars
  \let\commentend=\commentendchars}

\newcommand{\invisiblecomments}{%
  \let\onelinecomment=\empty
  \let\commentbegin=\empty
  \let\commentend=\empty}

\visiblecomments

\newlength{\blanklineskip}
\setlength{\blanklineskip}{0.66084ex}

\newcommand{\hsindent}[1]{\quad}% default is fixed indentation
\let\hspre\empty
\let\hspost\empty
\newcommand{\NB}{\textbf{NB}}
\newcommand{\Todo}[1]{$\langle$\textbf{To do:}~#1$\rangle$}

\EndFmtInput
\makeatother
%
%
%
%
%
%
% This package provides two environments suitable to take the place
% of hscode, called "plainhscode" and "arrayhscode". 
%
% The plain environment surrounds each code block by vertical space,
% and it uses \abovedisplayskip and \belowdisplayskip to get spacing
% similar to formulas. Note that if these dimensions are changed,
% the spacing around displayed math formulas changes as well.
% All code is indented using \leftskip.
%
% Changed 19.08.2004 to reflect changes in colorcode. Should work with
% CodeGroup.sty.
%
\ReadOnlyOnce{polycode.fmt}%
\makeatletter

\newcommand{\hsnewpar}[1]%
  {{\parskip=0pt\parindent=0pt\par\vskip #1\noindent}}

% can be used, for instance, to redefine the code size, by setting the
% command to \small or something alike
\newcommand{\hscodestyle}{}

% The command \sethscode can be used to switch the code formatting
% behaviour by mapping the hscode environment in the subst directive
% to a new LaTeX environment.

\newcommand{\sethscode}[1]%
  {\expandafter\let\expandafter\hscode\csname #1\endcsname
   \expandafter\let\expandafter\endhscode\csname end#1\endcsname}

% "compatibility" mode restores the non-polycode.fmt layout.

\newenvironment{compathscode}%
  {\par\noindent
   \advance\leftskip\mathindent
   \hscodestyle
   \let\\=\@normalcr
   \let\hspre\(\let\hspost\)%
   \pboxed}%
  {\endpboxed\)%
   \par\noindent
   \ignorespacesafterend}

\newcommand{\compaths}{\sethscode{compathscode}}

% "plain" mode is the proposed default.
% It should now work with \centering.
% This required some changes. The old version
% is still available for reference as oldplainhscode.

\newenvironment{plainhscode}%
  {\hsnewpar\abovedisplayskip
   \advance\leftskip\mathindent
   \hscodestyle
   \let\hspre\(\let\hspost\)%
   \pboxed}%
  {\endpboxed%
   \hsnewpar\belowdisplayskip
   \ignorespacesafterend}

\newenvironment{oldplainhscode}%
  {\hsnewpar\abovedisplayskip
   \advance\leftskip\mathindent
   \hscodestyle
   \let\\=\@normalcr
   \(\pboxed}%
  {\endpboxed\)%
   \hsnewpar\belowdisplayskip
   \ignorespacesafterend}

% Here, we make plainhscode the default environment.

\newcommand{\plainhs}{\sethscode{plainhscode}}
\newcommand{\oldplainhs}{\sethscode{oldplainhscode}}
\plainhs

% The arrayhscode is like plain, but makes use of polytable's
% parray environment which disallows page breaks in code blocks.

\newenvironment{arrayhscode}%
  {\hsnewpar\abovedisplayskip
   \advance\leftskip\mathindent
   \hscodestyle
   \let\\=\@normalcr
   \(\parray}%
  {\endparray\)%
   \hsnewpar\belowdisplayskip
   \ignorespacesafterend}

\newcommand{\arrayhs}{\sethscode{arrayhscode}}

% The mathhscode environment also makes use of polytable's parray 
% environment. It is supposed to be used only inside math mode 
% (I used it to typeset the type rules in my thesis).

\newenvironment{mathhscode}%
  {\parray}{\endparray}

\newcommand{\mathhs}{\sethscode{mathhscode}}

% texths is similar to mathhs, but works in text mode.

\newenvironment{texthscode}%
  {\(\parray}{\endparray\)}

\newcommand{\texths}{\sethscode{texthscode}}

% The framed environment places code in a framed box.

\def\codeframewidth{\arrayrulewidth}
\RequirePackage{calc}

\newenvironment{framedhscode}%
  {\parskip=\abovedisplayskip\par\noindent
   \hscodestyle
   \arrayrulewidth=\codeframewidth
   \tabular{@{}|p{\linewidth-2\arraycolsep-2\arrayrulewidth-2pt}|@{}}%
   \hline\framedhslinecorrect\\{-1.5ex}%
   \let\endoflinesave=\\
   \let\\=\@normalcr
   \(\pboxed}%
  {\endpboxed\)%
   \framedhslinecorrect\endoflinesave{.5ex}\hline
   \endtabular
   \parskip=\belowdisplayskip\par\noindent
   \ignorespacesafterend}

\newcommand{\framedhslinecorrect}[2]%
  {#1[#2]}

\newcommand{\framedhs}{\sethscode{framedhscode}}

% The inlinehscode environment is an experimental environment
% that can be used to typeset displayed code inline.

\newenvironment{inlinehscode}%
  {\(\def\column##1##2{}%
   \let\>\undefined\let\<\undefined\let\\\undefined
   \newcommand\>[1][]{}\newcommand\<[1][]{}\newcommand\\[1][]{}%
   \def\fromto##1##2##3{##3}%
   \def\nextline{}}{\) }%

\newcommand{\inlinehs}{\sethscode{inlinehscode}}

% The joincode environment is a separate environment that
% can be used to surround and thereby connect multiple code
% blocks.

\newenvironment{joincode}%
  {\let\orighscode=\hscode
   \let\origendhscode=\endhscode
   \def\endhscode{\def\hscode{\endgroup\def\@currenvir{hscode}\\}\begingroup}
   %\let\SaveRestoreHook=\empty
   %\let\ColumnHook=\empty
   %\let\resethooks=\empty
   \orighscode\def\hscode{\endgroup\def\@currenvir{hscode}}}%
  {\origendhscode
   \global\let\hscode=\orighscode
   \global\let\endhscode=\origendhscode}%

\makeatother
\EndFmtInput
%



%% ----------- Applications --------------
%% %format applyx a b = a "@" b
%% %format applyx a b = a "\," b

%% Function application with no arguments

%% Function application with 1 argument

%% Function application with 1 argument, without parenthesis

%% Function application with many arguments

%% ----------- End of Applications --------------


%% Sequences e1; ...; en; v

%% xs etc stole syntax used in examples.

%% Old version: | for BNF, [] for choice
%%%format choice = "[ \! ]"
%%%format `choice` = "\mathop{[ \! ]}"

%% New version: tall | for BNF, fat | for choice
%%format vbar = "\mathop{\bigm|}"
%%format choice = "\boldsymbol{|}"
%%format `choice` = "\mathop{\boldsymbol{|}}"




%% Arrays and tuples
%% %format array (x) = "\textbf{array}\{" x "\}"
%% %format arrKW = "\textbf{arr}"
%% %format arr (x) = arrKW "\{" x "\}"

%% %format defKW = "\textbf{def}"

% only used in comments, but still prettier this way:
%%%%format map = "\textbf{map}"
% format := = "\mapsto"
%  format effs x = "\!\!<\!\!" x "\!\!>\!\!" dots































%% ----------- Metatheory --------------

\section{Updateable references} \label{sec:store}

\begin{figure}
  $$
  \begin{array}{l}
    \textbf{Syntax extension} \\
    \begin{array}{lrcl}
      \text{References}         & r \\
      \text{Expressions}        & e   & ::= & \cdots \vb \ensuremath{\textbf{store}\;\Varid{s}\;\textbf{in}\,\{\Varid{e}\}} \\
      \text{Primops}            & op  & ::= & \cdots \vb \ensuremath{\textbf{alloc}} \vb \ensuremath{\textbf{read}} \vb \ensuremath{\textbf{write}} \\
      \text{Head values}        & hnf & ::= & \cdots \vb r \\
      \text{Execution contexts} & X   & ::= & \cdots \vb \ensuremath{\textbf{store}\;\Varid{s}\;\textbf{in}\,\{X\}} \\
      \text{Store}              & s   & ::= & \ensuremath{\Varid{r}\mapsto\!\Varid{v}} \vb \ensuremath{\Varid{s},\Varid{s}} \vb \ensuremath{\epsilon} \\
      \text{Store context}      & S   & ::= & \hole \vb \ensuremath{\Conid{S},\Varid{s}} \vb \ensuremath{\Varid{s},\Conid{S}} \\
    \end{array} \\
\\

  \textbf{Axiom extensions} \\
  \textit{Reference ops} \\
  \begin{array}{lrcll}
    \rulename{ref-alloc} & \ensuremath{\textbf{store}\;\Varid{s}\;\textbf{in}\,\{\textbf{alloc}(\Varid{v})\}} & \movesto & \ensuremath{\textbf{store}\;\Varid{s},\Varid{r}\mapsto\!\Varid{v}\;\textbf{in}\,\{\Varid{r}\}} & \text{\ensuremath{\Varid{r}} fresh} \\
    \rulename{ref-read} & \ensuremath{\textbf{store}\;\Conid{S} [\mskip1.5mu \Varid{r}\mapsto\!\Varid{v}\mskip1.5mu]\;\textbf{in}\,\{\textbf{read}(\Varid{r})\}} & \movesto & \ensuremath{\textbf{store}\;\Conid{S} [\mskip1.5mu \Varid{r}\mapsto\!\Varid{v}\mskip1.5mu]\;\textbf{in}\,\{\Varid{v}\}} \\
    \rulename{ref-write} &
        \ensuremath{\textbf{store}\;\Conid{S} [\mskip1.5mu \Varid{r}\mapsto\!\Varid{v}_{\mathrm{1}}\mskip1.5mu]\;\textbf{in}\,\{\textbf{write}\langle\Varid{r},\Varid{v}_{\mathrm{2}}\rangle\}}
        & \movesto & \ensuremath{\textbf{store}\;\Conid{S} [\mskip1.5mu \Varid{r}\mapsto\!\Varid{v}_{\mathrm{2}}\mskip1.5mu]\;\textbf{in}\,\{\langle\rangle\}} \\
  \end{array} \\
    \\
  \textit{Store propagation} \\
  \begin{array}{lrcll}
    % SEQ
    \rulename{st-seq-in} & \ensuremath{\textbf{store}\;\Varid{s}\;\textbf{in}\,\{\Varid{e}_{\mathrm{1}};\,\Varid{e}_{\mathrm{2}}\}}  & \movesto & \ensuremath{\textbf{store}\;\Varid{s}\;\textbf{in}\,\{\Varid{e}_{\mathrm{1}}\};\,\Varid{e}_{\mathrm{2}}} \\
    \rulename{st-seq-r} & \ensuremath{\textbf{store}\;\Varid{s}\;\textbf{in}\,\{\Varid{v}\};\,\Varid{e}}  & \movesto & \ensuremath{\Varid{v};\,\textbf{store}\;\Varid{s}\;\textbf{in}\,\{\Varid{e}\}} \\
    % UNIFY
    \rulename{st-unify-in} & \ensuremath{\textbf{store}\;\Varid{s}\;\textbf{in}\,\{\Varid{e}_{\mathrm{1}}\mathrel{=}\Varid{e}_{\mathrm{2}}\}}  & \movesto & \ensuremath{\textbf{store}\;\Varid{s}\;\textbf{in}\,\{\Varid{e}_{\mathrm{1}}\}\mathrel{=}\Varid{e}_{\mathrm{2}}} \\
    \rulename{st-unify-r} & \ensuremath{\textbf{store}\;\Varid{s}\;\textbf{in}\,\{\Varid{v}\}\mathrel{=}\Varid{e}}  & \movesto & \ensuremath{\Varid{v}\mathrel{=}\textbf{store}\;\Varid{s}\;\textbf{in}\,\{\Varid{e}\}} \\
    \rulename{st-unify} & \ensuremath{\Varid{v}_{\mathrm{1}}\mathrel{=}\textbf{store}\;\Varid{s}\;\textbf{in}\,\{\Varid{v}_{\mathrm{2}}\}}  & \movesto & \ensuremath{\Varid{v}_{\mathrm{1}}\mathrel{=}\Varid{v}_{\mathrm{2}};\,\textbf{store}\;\Varid{s}\;\textbf{in}\,\{\Varid{v}_{\mathrm{1}}\}} \\
    % CHOICE
    \rulename{st-choice} & \ensuremath{\textbf{store}\;\Varid{s}\;\textbf{in}\,\{\Varid{e}_{\mathrm{1}}\hspace{1mm}\mathop{\rule[-0.3mm]{0.5mm}{2.5mm}}\hspace{1mm}\Varid{e}_{\mathrm{2}}\}} & \movesto & \ensuremath{\textbf{store}\;\Varid{s}\;\textbf{in}\,\{\Varid{e}_{\mathrm{1}}\}\hspace{1mm}\mathop{\rule[-0.3mm]{0.5mm}{2.5mm}}\hspace{1mm}\Varid{e}_{\mathrm{2}}} \\
    % DEF
    \rulename{st-def-in} & \ensuremath{\textbf{store}\;\Varid{s}\;\textbf{in}\,\{\exists\Varid{x}.\,\Varid{e}\}} & \movesto & \ensuremath{\exists\Varid{x}.\,\textbf{store}\;\Varid{s}\;\textbf{in}\,\{\Varid{e}\}} & e \not\in v \\
    \rulename{st-def-out} & \ensuremath{\exists\Varid{x}.\,\textbf{store}\;\Varid{s}\;\textbf{in}\,\{\Varid{v}\}} & \movesto & \ensuremath{\textbf{store}\;\Varid{s}\;\textbf{in}\,\{\exists\Varid{x}.\,\Varid{v}\}} \\
    % SPLIT
    \rulename{st-split-value} & \ensuremath{\textbf{split}\{\textbf{store}\;\Varid{s}\;\textbf{in}\,\{\Varid{v}\},\Varid{f},\Varid{g}\}} & \movesto &
        \ensuremath{\textbf{store}\;\Varid{s}\;\textbf{in}\,\{\textbf{split}\{\Varid{v},\Varid{f},\Varid{g}\}\}} \\

    \rulename{st-split-choice} & \ensuremath{\textbf{split}\{\textbf{store}\;\Varid{s}\;\textbf{in}\,\{\Varid{v}\hspace{1mm}\mathop{\rule[-0.3mm]{0.5mm}{2.5mm}}\hspace{1mm}\Varid{e}\},\Varid{f},\Varid{g}\}} & \movesto &
      \ensuremath{\textbf{store}\;\Varid{s}\;\textbf{in}\,\{\textbf{split}\{\Varid{v}\hspace{1mm}\mathop{\rule[-0.3mm]{0.5mm}{2.5mm}}\hspace{1mm}\Varid{e},\Varid{f},\Varid{g}\}\}} \\

    % FAIL
    %% \rulename{st-wrong} & |store s wrong| & \movesto & |wrong| \\
  \end{array} \\
 \\
  \textit{Store copying} \\
  \begin{array}{lrcll}
    % SPLIT
    \rulename{st-split-copy} &
      \begin{array}[t]{@{}l}
        \ensuremath{\textbf{store}\;\Varid{s}\;\textbf{in}\\\{\textbf{split}\{\Varid{e},\Varid{f},\Varid{g}\}\}}
      \end{array} & \movesto &
      \begin{array}[t]{@{}l}
        \ensuremath{\textbf{store}\;\Varid{s}\;\textbf{in}\\\{\textbf{split}\{\textbf{store}\;\Varid{s}\;\textbf{in}\,\{\Varid{e}\},\Varid{f},\Varid{g}\}\}} \\
      \end{array} & \text{\ensuremath{\Varid{canCopy}\;\Varid{e}}} \\
    % FAIL
    \rulename{st-fail} & \ensuremath{\textbf{store}\;\Varid{s}\;\textbf{in}\,\{\textbf{fail}\}} & \movesto & \ensuremath{\textbf{fail}} \\
    % STORE
    \rulename{st-store} & \ensuremath{\textbf{store}\;\Varid{s}_{\mathrm{1}}\;\textbf{in}\,\{\textbf{store}\;\Varid{s}_{\mathrm{2}}\;\textbf{in}\,\{\Varid{v}\}\}} & \movesto & \ensuremath{\textbf{store}\;\Varid{s}_{\mathrm{2}}\;\textbf{in}\,\{\Varid{v}\}} \\
  \end{array} \\
    \\
    \textit{Unification} \\
    \begin{array}{lrcll}
      \text{Extension with the obvious axioms making equal references unify,
        and anything else fail.}
    \end{array} \\
  \\
  \textit{Top level} \\
  \begin{array}{lrcll}
    \text{Start top level reduction of \ensuremath{\Varid{e}} with \ensuremath{\textbf{store}\;\epsilon\;\textbf{in}\,\{\Varid{e}\}}.}
  \end{array} \\

    \\
    \text{\ensuremath{\Varid{canCopy}\;\Varid{e}} is true if \ensuremath{\Varid{e}} contains no \ensuremath{\textbf{store}}, \ensuremath{\Varid{e}} isn't \ensuremath{\Varid{v}}, \ensuremath{\Varid{e}} isn't \ensuremath{\Varid{v}\hspace{1mm}\mathop{\rule[-0.3mm]{0.5mm}{2.5mm}}\hspace{1mm}\Varid{e}_{\mathrm{1}}}, nor \ensuremath{\textbf{fail}}}

  \end{array}
   $$
  \caption{The Verse calculus: store axioms} \label{fig:store-axioms}
\end{figure}


To add updateable references we extend the system with syntax and rules from figure~\ref{fig:store-axioms}.
The \ensuremath{\textbf{store}\;\Varid{s}\;\textbf{in}\,\{\Varid{e}\}} indicates that \ensuremath{\Varid{e}} should be reduced using the store \ensuremath{\Varid{s}}.
A store is simply a mapping from references to values (one mapping being \ensuremath{\Varid{r}\mapsto\!\Varid{v}}).
A reference is some opaque type that supports equality (unification) and creation of a new reference.

There are three new operations, \ensuremath{\textbf{alloc}}, \ensuremath{\textbf{read}}, and \ensuremath{\textbf{write}}.
Their interaction with the store can be seen from the axioms.
The \ensuremath{\textbf{alloc}(\Varid{v})} operation creates a new reference and adds a binding with \ensuremath{\Varid{v}} to the store.
The \ensuremath{\textbf{read}(\Varid{r})} operation retrieves the value for reference \ensuremath{\Varid{r}} from the store,
and \ensuremath{\textbf{write}\langle\Varid{r},\Varid{v}\rangle} updates the reference \ensuremath{\Varid{r}} with \ensuremath{\Varid{v}} in the store.
Note that these reductions only happen when the store is exactly at the operation.

Most of the store axioms just move the store around from left-to-right, ``visiting'' every subexpression,
and moving on when the subexpression has reduced to a value.
This way every subexpression (in a non-looping program) will get a ``visit'' and all reference operations
can reduce.

The interesting rules are those that cover \ensuremath{\textbf{split}}.
When reducing the \ensuremath{\Varid{e}} in \ensuremath{\textbf{split}\{\Varid{e},\Varid{f},\Varid{g}\}} the reduction of \ensuremath{\Varid{e}} is treated as a transaction, {\ie},
if the reduction fails any changes to the store are rolled back.
This is handled by a combination of \rulename{st-split-copy}, \rulename{st-fail}, and \rulename{st-store}.
First, when a \ensuremath{\textbf{store}} is at a \ensuremath{\textbf{split}\{\Varid{e},\Varid{f},\Varid{g}\}} then the store is copied to reduce the \ensuremath{\Varid{e}}.
If the evaluation of \ensuremath{\Varid{e}} fails, then rule \rulename{st-fail} will simply abandon the copy of the store.
If the evaluation of \ensuremath{\Varid{e}} is successful then the \rulename{st-value} or \rulename{st-choice} will move the store out of the \ensuremath{\textbf{split}}.
In that position the \rulename{st-store} rule can ``commit'' and throw away the old store, replacing it with the new.
Note that the \rulename{st-split-copy} has a side condition to stop it from firing when it has just fired,
and also when the \ensuremath{\Varid{e}} is a result.

To support limited store operations ({\eg}, \ensuremath{\textbf{read}}, but not \ensuremath{\textbf{write}}) we can equip the store with a set of currently allowed operations.
We also need some extra expressions that modify this set as the store ``moves'' around.

\lennart{These rewrite rules assume that choices/\ensuremath{\textbf{fail}} are always inside a \ensuremath{\textbf{split}}, otherwise the store can get lost.}
